% References.
\chapter*{References}\addcontentsline{toc}{chapter}{References}
\markboth{References}{References}

\begin{description}[leftmargin=0pt,labelindent=0pt,itemsep=3pt]
\item Baarda, W. (1968). \emph{A testing procedure for use in geodetic networks.} Netherlands
  Geodetic Commission, Publications on Geodesy, 2(5).
\item Ellmer, M. (2018). \emph{Contributions to GRACE gravity field recovery: improvements in
  dynamic orbit integration, stochastic modelling of the antenna offset correction, and
  co-estimation of satellite orientations.} Monographic Series TU Graz, Graz University of
  Technology.
\item Förstner, W. (1979). Ein Verfahren zur Schätzung von Varianz- und Kovarianzkomponenten.
  \emph{Allgemeine Vermessungs-Nachrichten} 86, 446--453.
\item Hutchinson, M. F. (1990). A stochastic estimator of the trace of the influence matrix for
  Laplacian smoothing splines. \emph{Communications in Statistics -- Simulation and Computation}
  19(2), 433--450.
\item Koch, K.-R. (1999). \emph{Parameter estimation and hypothesis testing in linear models.}
  2nd ed., Springer.
\item Koch, K.-R., Kusche, J. (2002). Regularization of geopotential determination from satellite
  data by variance components. \emph{Journal of Geodesy} 76, 259--268.
\item Kvas, A. (2019). \emph{Gravity field recovery from GRACE-type observations: daily solutions
  and process noise.} PhD thesis, Graz University of Technology.
\item Ljung, G. M., Box, G. E. P. (1978). On a measure of lack of fit in time series models.
  \emph{Biometrika} 65(2), 297--303.
\item Pope, A. J. (1976). \emph{The statistics of residuals and the detection of outliers.} NOAA
  Technical Report NOS 65 NGS 1.
\item Yang, Y. (1994). Robust estimation for dependent observations. \emph{Manuscripta Geodaetica}
  19, 10--17.
\end{description}
